\documentclass[11pt]{article}

\usepackage[margin=1in]{geometry}
\usepackage{amsmath, amssymb, amsthm}
\usepackage{natbib}
\usepackage{hyperref}
\usepackage{enumitem}

\setlength{\parskip}{4pt}
\setlength{\parindent}{0pt}

\title{Causal survival estimands under the AFT data-fusion model\\
with HDP error}
\author{Working notes}
\date{\today}

\newcommand{\X}{X}
\newcommand{\A}{A}
\newcommand{\Sx}{S}
\newcommand{\tauf}{\tau}
\newcommand{\cf}{c}
\newcommand{\gf}{g}
\newcommand{\muf}{\mu}
\newcommand{\mzero}{m_{0}}
\newcommand{\Spot}{S}
\newcommand{\RCT}{\mathrm{RCT}}
\newcommand{\OS}{\mathrm{OS}}
\newcommand{\RMST}{\mathrm{RMST}}
\newcommand{\SD}{\mathrm{SD}}
\newcommand{\AF}{\mathrm{AF}}

\begin{document}
\maketitle

\section{Setup}

We work under the AFT decomposition
\begin{equation}
  \log T_{i} = \mzero(\X_{i}, \Sx_{i})
            + \A_{i}\,\tauf(\X_{i})
            + (1 - \Sx_{i})\,\A_{i}\,\cf(\X_{i})
            + \sigma\,\varepsilon_{i},
  \qquad \varepsilon_{i} \sim P,
  \label{eq:model}
\end{equation}
with $P$ assigned a hierarchical Dirichlet process (HDP) prior with Gaussian
kernels and atoms shared across the two data sources. Conditional on the
sampler state, the implied error CDF in source $s$ is
\begin{equation}
  \Phi_{s}(z) = \sum_{k=1}^{K} \pi_{s, k}\,
                \Phi^{\mathrm{N}}\!\left(\frac{z - \mu_{k}}{\tau_{k}}\right),
  \label{eq:hdp-cdf}
\end{equation}
where $\Phi^{\mathrm{N}}$ is the standard normal CDF,
$\{(\mu_{k}, \tau_{k})\}_{k=1}^{K}$ are the cluster atoms (shared across
sources via the top-level DP), and $\pi_{s,k}$ are the source-specific
weights.

The treatment-specific AFT shift at target population $s_{t}$ is
\begin{equation}
  \eta_{a}(x, s_{t}) := \mzero(x, s_{t})
                    + a\,\tauf(x)
                    + (1 - s_{t})\,a\,\cf(x),
  \qquad a \in \{0, 1\}.
\end{equation}
The potential survival function is
\begin{equation}
  \Spot_{a}(t \mid x, s_{t})
  = 1 - \Phi_{s_{t}}\!\left(\frac{\log t - \eta_{a}(x, s_{t})}{\sigma}\right),
  \label{eq:potential-S}
\end{equation}
and we collect causal estimands as functionals of
$\{\Spot_{0}(\cdot \mid x, s_{t}), \Spot_{1}(\cdot \mid x, s_{t})\}$.

For the closed-form expressions below it is convenient to define, for each
component $k$ of the HDP mixture,
\begin{equation}
  \xi_{a, k}(x, s_{t}) := \eta_{a}(x, s_{t}) + \sigma\mu_{k},
  \qquad
  \omega_{k} := \sigma\tau_{k},
\end{equation}
so that the per-component survival function under the Gaussian-kernel HDP
is lognormal:
$1 - \Phi^{\mathrm{N}}\bigl((\log t - \xi_{a, k})/\omega_{k}\bigr)$.

\section{The survival difference $\Delta_{\SD}$}

\subsection{Definition}

The survival (risk) difference at time $t$ is
\begin{equation}
  \Delta_{\SD}(t; x, s_{t})
  := \Spot_{1}(t \mid x, s_{t}) - \Spot_{0}(t \mid x, s_{t}).
\end{equation}

\subsection{Closed form under the HDP}

\begin{equation}
  \boxed{\;
    \Delta_{\SD}(t; x, s_{t})
    = \sum_{k=1}^{K} \pi_{s_{t}, k}\!\left[
        \Phi^{\mathrm{N}}\!\left(\frac{\log t - \xi_{0, k}(x, s_{t})}{\omega_{k}}\right)
      - \Phi^{\mathrm{N}}\!\left(\frac{\log t - \xi_{1, k}(x, s_{t})}{\omega_{k}}\right)
    \right].
  \;}
  \label{eq:SD-closed}
\end{equation}

\subsection{Interpretation}

$\Delta_{\SD}(t; x, s_{t})$ is the absolute increase in the probability of
surviving past time $t$ attributable to assigning treatment versus control
for individuals with covariates $x$ in target population $s_{t}$. It lives
on the probability scale, is bounded between $-1$ and $1$, and has a direct
clinical reading: ``treatment increases the chance of surviving to year $t$
by $\Delta_{\SD}$ percentage points.''

This is the absolute risk reduction (ARR), or its survival-scale analogue.
It is the estimand most directly tied to clinical decision-making at a fixed
horizon and is the natural quantity for computing the number-needed-to-treat
$1/|\Delta_{\SD}(t; x, s_{t})|$ when the contrast is favourable
\citep{Cook1995}.

\paragraph{When to use.} When a single clinically meaningful horizon
exists --- for example, ``5-year survival in early-stage cancer'' or
``1-year survival post-transplant'' --- the SD is the most interpretable
summary. It depends on the error distribution $\Phi_{s_{t}}$ through
\eqref{eq:SD-closed} and is therefore sensitive to misspecification of the
tail when $t$ extends beyond the RCT follow-up; in the present framework
that sensitivity is absorbed by the HDP error.

\paragraph{Limitations.} $\Delta_{\SD}(t; x, s_{t})$ is a pointwise contrast
at a single $t$ and can give a misleading picture when survival curves
cross or when the gain is concentrated outside the chosen window. It does
not summarise the full survival benefit. For these reasons, regulatory and
HTA guidance recommends pairing the SD with an integrated summary such as
the RMST \citep{Royston2013}.

\section{The restricted mean survival time difference $\Delta_{\RMST}$}

\subsection{Definition}

The RMST under treatment $a$ at horizon $t^{*}$ is
\begin{equation}
  \mu_{a}(t^{*}; x, s_{t}) := \int_{0}^{t^{*}}\!
  \Spot_{a}(u \mid x, s_{t})\,du,
\end{equation}
and the RMST contrast is
\begin{equation}
  \Delta_{\RMST}(t^{*}; x, s_{t}) := \mu_{1}(t^{*}; x, s_{t}) - \mu_{0}(t^{*}; x, s_{t}).
\end{equation}

\subsection{Closed form under the HDP}

The per-component lognormal RMST has the standard expression
\begin{equation}
  \mu_{a}^{(k)}(t^{*}; x, s_{t})
  = e^{\,\xi_{a, k} + \omega_{k}^{2}/2}\,
    \Phi^{\mathrm{N}}\!\left(\frac{\log t^{*} - \xi_{a, k} - \omega_{k}^{2}}{\omega_{k}}\right)
  + t^{*}\!\left[1 - \Phi^{\mathrm{N}}\!\left(\frac{\log t^{*} - \xi_{a, k}}{\omega_{k}}\right)\right],
\end{equation}
and the contrast is the source-specific mixture
\begin{equation}
  \boxed{\;
    \Delta_{\RMST}(t^{*}; x, s_{t})
    = \sum_{k=1}^{K} \pi_{s_{t}, k}\bigl[
        \mu_{1}^{(k)}(t^{*}; x, s_{t}) - \mu_{0}^{(k)}(t^{*}; x, s_{t})
      \bigr].
  \;}
  \label{eq:RMST-closed}
\end{equation}

\subsection{Interpretation}

$\Delta_{\RMST}(t^{*}; x, s_{t})$ is the expected gain in survival time
within the window $[0, t^{*}]$ attributable to treatment, for individuals
with covariates $x$ in target population $s_{t}$. It lives on the
\emph{time scale}, in the same units as the outcome, and has a direct
reading: ``treatment yields $\Delta_{\RMST}$ additional units of expected
life within the first $t^{*}$ units.''

The RMST has been advanced as a principal alternative to the hazard ratio,
particularly when proportional hazards is implausible
\citep{Royston2011, Royston2013, Uno2014}. Its appeal is threefold: it has
a transparent units-of-time interpretation; it is well-defined under
non-proportional hazards including crossing survival curves; and it
integrates the full survival benefit within the window rather than
reporting a contrast at a single point. \citet{Uno2014} review its use in
oncology trials and provide guidance on the choice of $t^{*}$.

\paragraph{When to use.} The RMST is the natural headline estimand under
the long-term-extrapolation goal of this work. Reporting
$\Delta_{\RMST}(t^{*}; x, s_{t})$ at multiple horizons --- e.g.\ at
$t^{*}$ inside the trial window, at $t^{*}$ just past it, and at
$t^{*}$ far beyond it --- traces how the survival benefit evolves with
horizon and exposes the contribution of the HDP-tail to the inference.

\paragraph{Choice of $t^{*}$.} The horizon $t^{*}$ should be chosen
ex ante and motivated by clinical relevance or by regulatory context.
\citet{Royston2013} recommend $t^{*}$ at or near the largest follow-up
time at which a non-negligible proportion of the at-risk set remains
under observation. In the present setting that recommendation is too
conservative: the central methodological point is to extrapolate beyond
the RCT follow-up, so $t^{*}$ should deliberately exceed it.

\paragraph{Limitations.} The RMST depends on $t^{*}$, and substantive
conclusions can shift when $t^{*}$ is varied. Beyond the OS follow-up,
$\Delta_{\RMST}$ inherits the full extrapolation uncertainty of the
posterior and should be reported with appropriate credible intervals.

\section{The acceleration factor $\AF$}

\subsection{Definition}

Under the AFT model \eqref{eq:model}, the contrast between treatment and
control at fixed covariates $x$ is a constant additive shift on the
log-time scale. Equivalently, on the time scale it is a multiplicative
\emph{acceleration factor}. The causal acceleration factor at $x$ is
\begin{equation}
  \AF(x) := \exp\bigl\{\eta_{1}(x, s_{t}) - \eta_{0}(x, s_{t})\bigr\}
         = \exp\bigl\{\tauf(x) + (1 - s_{t})\,\cf(x)\bigr\}.
\end{equation}
For the causal contrast in either population, we use $s_{t} = 1$ (where
\eqref{eq:model} is unconfounded by construction) and the confounding
term vanishes:
\begin{equation}
  \boxed{\;
    \AF_{\mathrm{causal}}(x) = \exp\{\tauf(x)\}.
  \;}
  \label{eq:AF-causal}
\end{equation}
By transportability of potential outcomes (Assumption~4 of the framework
document), $\AF_{\mathrm{causal}}(x)$ is the same in the RCT and OS
populations: it is a property of the causal mechanism, not of the source.

The corresponding observed contrast in the OS,
$\AF_{\mathrm{obs, OS}}(x) = \exp\{\tauf(x) + \cf(x)\}$, is the
confounded analogue. Its discrepancy from $\AF_{\mathrm{causal}}(x)$ is
precisely $\exp\{\cf(x)\}$ and quantifies the multiplicative bias on the
time scale due to unmeasured confounding in the OS.

\subsection{Closed form and HDP-independence}

The acceleration factor is a direct functional of the BART forest for
$\tauf$ and does not involve the error distribution. In particular,
\begin{equation}
  \AF_{\mathrm{causal}}(x) = \exp\{\tauf(x)\}
\end{equation}
is computable at each MCMC iteration as $\exp\{\tauf^{(r)}(x)\}$, without
reference to the cluster atoms or weights of the HDP. The full
posterior of $\AF_{\mathrm{causal}}(x)$ is therefore available at the
same cost as the posterior of $\tauf(x)$.

\subsection{Interpretation}

$\AF_{\mathrm{causal}}(x)$ is the ratio by which treatment multiplies an
individual's time scale: a unit with covariates $x$ assigned treatment
experiences event times that are $\AF_{\mathrm{causal}}(x)$ times those
under control. The reading is symmetric: $\AF_{\mathrm{causal}}(x) = 1.5$
means treatment extends time-to-event by 50\%; $\AF_{\mathrm{causal}}(x) =
0.8$ means treatment shortens it by 20\%.

The acceleration-factor interpretation is the defining feature of AFT
models and is articulated in detail in \citet{Wei1992},
\citet{KalbfleischPrentice2002}, and \citet{CoxOakes1984}. Compared with
the hazard ratio, the AF has the advantages that (i) it is directly
interpretable on the time scale rather than on the hazard scale, (ii) it
is meaningful even when hazards are non-proportional, and (iii) under
AFT it equals the ratio of median (and any other quantile) survival
times:
\begin{equation}
  \frac{\mathrm{Quantile}_{q}\bigl(T(1) \mid x\bigr)}
       {\mathrm{Quantile}_{q}\bigl(T(0) \mid x\bigr)}
  = \AF_{\mathrm{causal}}(x)
  \qquad \text{for all } q \in (0, 1),
\end{equation}
a property sometimes called \emph{quantile equivariance}
\citep{Wei1992, Hougaard1999}.

\paragraph{When to use.} The AF is the natural estimand whenever the goal
is to summarise the treatment effect by \emph{a single number per
covariate value}, free of horizon, free of error-distribution assumptions,
and on a transparent scale. In data-fusion settings, this is particularly
convenient: $\AF_{\mathrm{causal}}(x)$ depends only on $\tauf(x)$ and is
therefore unaffected by the extrapolation uncertainty that enters the
RMST and SD via the HDP tail. It can be reported alongside the SD and
RMST as a horizon-free summary.

\paragraph{Limitations.} The AF carries the full weight of the AFT
structural assumption: that the treatment effect is a single time-scale
acceleration that does not vary with $t$. When the true effect waxes or
wanes over time --- as for some immune-mediated therapies --- a single
$\AF_{\mathrm{causal}}(x)$ misrepresents the effect. In such cases,
SD and RMST at multiple $t$ or $t^{*}$ are more informative.

A second caveat is that the AF is a multiplicative-on-time-scale measure
and does not by itself communicate absolute benefit. A 50\% acceleration
on a baseline median of one month and on a baseline median of ten years
have very different clinical implications. Pairing the AF with the SD
or RMST closes this gap.

\section{Population-averaged versions}

Each of the three estimands has a population-averaged counterpart obtained
by marginalising over the covariate distribution of the target population:
\begin{align}
  \Delta_{\SD}(t; s_{t})
  &= \mathbb{E}\bigl[\Delta_{\SD}(t; \X, s_{t}) \mid \Sx = s_{t}\bigr], \\
  \Delta_{\RMST}(t^{*}; s_{t})
  &= \mathbb{E}\bigl[\Delta_{\RMST}(t^{*}; \X, s_{t}) \mid \Sx = s_{t}\bigr], \\
  \AF_{\mathrm{causal}}
  &= \mathbb{E}\bigl[\AF_{\mathrm{causal}}(\X)\bigr]
   \quad \text{or} \quad
   \exp\bigl(\mathbb{E}[\tauf(\X)]\bigr),
\end{align}
where the last definition distinguishes the population-averaged AF (which
preserves the multiplicative interpretation pointwise but is no longer a
ratio of population means) from the AF of the population-averaged log-time
shift (which is). The two coincide only under degenerate $\tauf(\X)$.
Pointwise reporting via the posterior projection summaries is generally
preferable to a single marginal number unless heterogeneity is small.

The Monte Carlo estimator is in each case the empirical average over
evaluation rows of the corresponding posterior draw, with Bayesian-bootstrap
weights if covariate-distribution uncertainty is to be propagated.

\section{Summary}

\begin{itemize}[leftmargin=*]
  \item $\Delta_{\SD}(t; x, s_{t})$ --- absolute survival-probability
        gain at fixed $t$. Probability scale; depends on HDP tail.
        Natural for fixed clinical horizons and number-needed-to-treat.
  \item $\Delta_{\RMST}(t^{*}; x, s_{t})$ --- expected life-time gain
        within $[0, t^{*}]$. Time scale; depends on HDP tail.
        Natural headline estimand under the long-term extrapolation goal.
  \item $\AF_{\mathrm{causal}}(x) = \exp\{\tauf(x)\}$ --- multiplicative
        time-scale acceleration. Horizon-free; HDP-free; assumes
        the AFT shift structure. Natural complement summarising treatment
        heterogeneity in a single covariate-specific number.
\end{itemize}

\begin{thebibliography}{99}

\bibitem[Cook and Sackett(1995)]{Cook1995}
Cook, R.~J.\ and Sackett, D.~L. (1995).
\newblock The number needed to treat: a clinically useful measure of
treatment effect.
\newblock \emph{BMJ}, 310(6977):452--454.

\bibitem[Cox and Oakes(1984)]{CoxOakes1984}
Cox, D.~R.\ and Oakes, D. (1984).
\newblock \emph{Analysis of Survival Data}.
\newblock Chapman \& Hall, London.

\bibitem[Hougaard(1999)]{Hougaard1999}
Hougaard, P. (1999).
\newblock Fundamentals of survival data.
\newblock \emph{Biometrics}, 55(1):13--22.

\bibitem[Kalbfleisch and Prentice(2002)]{KalbfleischPrentice2002}
Kalbfleisch, J.~D.\ and Prentice, R.~L. (2002).
\newblock \emph{The Statistical Analysis of Failure Time Data}, 2nd ed.
\newblock Wiley, New York.

\bibitem[Royston and Parmar(2011)]{Royston2011}
Royston, P.\ and Parmar, M.~K.~B. (2011).
\newblock The use of restricted mean survival time to estimate the
treatment effect in randomized clinical trials when the proportional
hazards assumption is in doubt.
\newblock \emph{Statistics in Medicine}, 30(19):2409--2421.

\bibitem[Royston and Parmar(2013)]{Royston2013}
Royston, P.\ and Parmar, M.~K.~B. (2013).
\newblock Restricted mean survival time: an alternative to the hazard
ratio for the design and analysis of randomized trials with a
time-to-event outcome.
\newblock \emph{BMC Medical Research Methodology}, 13:152.

\bibitem[Uno et al.(2014)]{Uno2014}
Uno, H., Claggett, B., Tian, L., Inoue, E., Gallo, P., Miyata, T.,
Schrag, D., Takeuchi, M., Uyama, Y., Zhao, L., Skali, H., Solomon, S.,
Jacobus, S., Hughes, M., Packer, M., and Wei, L.~J. (2014).
\newblock Moving beyond the hazard ratio in quantifying the between-group
difference in survival analysis.
\newblock \emph{Journal of Clinical Oncology}, 32(22):2380--2385.

\bibitem[Wei(1992)]{Wei1992}
Wei, L.~J. (1992).
\newblock The accelerated failure time model: a useful alternative to the
Cox regression model in survival analysis.
\newblock \emph{Statistics in Medicine}, 11(14--15):1871--1879.

\end{thebibliography}

\end{document}